El estado del arte se organiza en \textbf{seis ejes} (E1--E6) del Anexo~B. Cada eje responde a una
pregunta sobre OpenAPI, \textit{matching} estructural o semántico, arquitectura de servicios o
delimitación del foco. La \Cref{tab:sintesis-ejes} resume hallazgos convergentes; el Anexo~B detalla
cada paper.

\subsection{Eje E1 --- OpenAPI como estándar}

\textbf{Papers:} Casas et al.\ (2021), Mainas et al.\ (2023), Santos y Casas (2025).

Casas et al.\ (2021) realizan un mapeo sistemático de 47 artículos primarios (2011--2020) sobre
usos de OpenAPI/Swagger. Confirman que es el \textbf{estándar de facto} para describir APIs REST de
forma \textit{machine-readable}: el 68\,\% de los trabajos mapeados propone herramientas o
automatización; el 43\,\% mejora documentación. Sin embargo, ningún estudio del eje vincula esos
contratos con modelos de referencia de negocio bancario tipo BIAN.

Mainas et al.\ (2023) avanzan hacia \textbf{OpenAPI enriquecido semánticamente}. Proponen
anotaciones (\texttt{x-refersTo}, \texttt{x-kindOf}, \texttt{x-mapsTo}, etc.), ontología REST/Hydra,
traducción OpenAPI $\rightarrow$ RDF y consultas SPARQL sobre un corpus de más de 300 contratos de
SwaggerHub. Aportan \textbf{técnicas de representación y normalización} candidatas para calcular
\SemanticScore{}; no definen S ni \AlignmentScore{}, y operan en dominio genérico sin
\ServiceDomain{} BIAN.

Santos y Casas (2025), mediante Design Science Research, construyen un catálogo de 59 métricas de
\textbf{API Gateway} (latencia, tráfico, errores) evaluadas en 68 productos comerciales. Delimitan un
borde útil para este estudio: la calidad operativa en producción \textbf{no equivale} a la coherencia
documental OpenAPI $\leftrightarrow$ modelo de negocio.

\textbf{Conclusión E1:} OpenAPI está maduro como artefacto técnico; la literatura no resuelve su
\textbf{alineación} con capacidades de negocio estandarizadas en banca.

\subsection{Eje E2 --- Matching estructural}

\textbf{Paper:} Shvaiko y Euzenat (2005).

Shvaiko y Euzenat (2005) sistematizan el \textbf{\textit{schema matching}} en una taxonomía que
cruza aproximación exacta/approximada, evidencia sintáctica/semántica/externa y nivel
elemento/estructura. El marco es anterior a OpenAPI y JSON Schema ---no aborda contratos REST ni
dominios bancarios---, pero fundamenta el \StructuralScore{} al contrastar \textit{paths},
\textit{operations} y \textit{schemas} con la organización de un \ServiceDomain{} BIAN
(\Behavior{}, \BusinessObject{}) como problema de correspondencia entre esquemas.

\textbf{Conclusión E2:} Existe base teórica sólida para correspondencias estructurales; falta
aplicarla a \textbf{pares OpenAPI bancario $\leftrightarrow$ BIAN} con métrica reproducible.

\subsection{Eje E3 --- Matching semántico y ontologías}

\textbf{Papers:} El Bayadh et al.\ (2015), Chandrasekaran y Mago (2022), Santos Sousa et al.\ (2026).

Tres líneas convergen. El Bayadh et al.\ (2015) clasifican \textbf{alineación ontológica} en medidas
terminológicas, estructurales, extensionales y semánticas; advierten que ninguna medida aislada basta.
Chandrasekaran y Mago (2022) trazan la evolución de \textbf{similitud semántica}
desde recursos léxicos hasta \textit{embeddings} y \textit{transformers}. Santos Sousa
et al.\ (2026) revisan \textbf{\textit{embeddings} para ontology matching}, con limitaciones de
supervisión, escalabilidad y correspondencias complejas.

Ninguno aborda OpenAPI ni BIAN directamente; todos aportan técnicas candidatas para
\SemanticScore{} y, en parte, cobertura conceptual (\CoverageScore{}).

\textbf{Conclusión E3:} La similitud semántica está bien documentada; no existe propuesta integrada
que la aplique a \textbf{vocabulario OpenAPI bancario vs.\ conceptos BIAN} dentro de un modelo de
alineación único.

\subsection{Eje E4 --- Microservicios e interfaces}

\textbf{Paper:} Schmidt y Thiry (2020).

Schmidt y Thiry (2020) revisan 27 estudios primarios (de 715 registros, 2013--2019) sobre
estrategias de \textbf{identificación de microservicios} con MDE y DDD. La síntesis muestra cómo la
descomposición desde modelos de dominio hacia límites de servicio contextualiza por qué los
\ServiceDomain{} BIAN son referencia de negocio y las APIs su capa de exposición. Pocos trabajos
integran MDE y DDD conjuntamente; ninguno mide coherencia entre contrato OpenAPI y modelo BIAN.

\textbf{Conclusión E4:} La literatura arquitectónica relaciona dominio e interfaces; \textbf{no
cuantifica} desalineación OpenAPI--BIAN.

\subsection{Eje E5 --- Testing y calidad API (borde del estudio)}

\textbf{Papers:} Ehsan et al.\ (2022), Golmohammadi et al.\ (2024), Otieno y Araka (2026); Tanveer
et al.\ (2025), Vitorino et al.\ (2025).

Cinco revisiones sistemáticas cubren \textbf{testing REST}, ML para pruebas, \textbf{vulnerabilidades}
y generación adversarial de datos. En todos los casos OpenAPI entra como \textbf{entrada} a pruebas
automatizadas, no como artefacto a contrastar con un \ServiceDomain{} BIAN.

Ehsan et al.\ (2022) incluyen 16 artículos sobre metodologías de testing RESTful. La
\textit{code coverage} es la métrica principal de efectividad.

Golmohammadi et al.\ (2024) analizan 92 papers con taxonomía de técnicas (\textit{fuzzing},
\textit{search-based}, model-based). La mayoría consume contratos OpenAPI para generar casos de
prueba sobre endpoints ya desplegados.

Otieno y Araka (2026) revisan ML aplicado a testing RESTful; reportan datasets pequeños y costo de
etiquetado como limitaciones recurrentes.

Tanveer et al.\ (2025) sistematizan detección de vulnerabilidades en APIs REST; el foco es seguridad
\textit{runtime}, ajeno al \AlignmentScore{} documental.

Vitorino et al.\ (2025) abordan aprendizaje adversarial para pruebas automatizadas. Confirman que la
literatura de calidad operativa \textbf{no sustituye} un procedimiento de alineación OpenAPI--BIAN.

\textbf{Conclusión E5:} Confirma la delimitación del plan (\Cref{sec:alcances}): testing, seguridad
y QoS operativa \textbf{quedan fuera}; el producto central es el \AlignmentScore{} documental.

\subsection{Eje E6 --- Selección y recomendación de servicios}

\textbf{Paper:} Mahanra Rao et al.\ (2025).

Mahanra Rao et al.\ (2025) revisan 25 estudios (2020--2024, PRISMA) sobre enfoques
\textbf{BERT/transformers} para selección y recomendación de servicios web con foco en predicción de
QoS. La representación semántica de descripciones de servicio es adyacente a S, pero el objetivo es
\textbf{ranking/consumo}, no validar un contrato frente a un \ServiceDomain{} de referencia. Los
autores señalan alto costo computacional, datasets limitados e interpretabilidad reducida.

\textbf{Conclusión E6:} Contexto semántico de servicios; \textbf{no sustituye} un modelo de
alineación OpenAPI $\leftrightarrow$ BIAN.

\begin{table}[H]
  \centering
  \caption{Síntesis del estado del arte por ejes (E1--E6)}
  \label{tab:sintesis-ejes}
  \small
  \begin{tabular}{@{}p{0.06\textwidth}p{0.20\textwidth}p{0.22\textwidth}p{0.16\textwidth}p{0.24\textwidth}@{}}
    \toprule
    \textbf{Eje} & \textbf{Papers} & \textbf{Hallazgo principal} & \textbf{Aporte S/E/C} & \textbf{Vacío OpenAPI--BIAN} \\
    \midrule
    \textbf{E1} &
    Casas 2021; Mainas 2023; Santos y Casas 2025 &
    OpenAPI consolidado; enriquecimiento semántico posible &
    S (Mainas); contexto estándar &
    Sin vínculo con \ServiceDomain{} BIAN \\
    \textbf{E2} &
    Shvaiko 2005 &
    Taxonomía de \textit{schema matching} &
    E &
    Sin aplicación REST/banca/BIAN \\
    \textbf{E3} &
    El Bayadh 2015; Chandrasekaran 2022; Santos Sousa 2026 &
    Ontología matching + embeddings + STS &
    S; C parcial &
    Sin contratos OpenAPI ni dominio bancario \\
    \textbf{E4} &
    Schmidt 2020 &
    Dominio (DDD/MDE) $\rightarrow$ microservicios &
    Contexto dominio$\rightarrow$interfaz (no \StructuralScore{} operativo) &
    Sin métrica de alineación \\
    \textbf{E5} &
    Ehsan 2022; Golmohammadi 2024; Otieno 2026; Tanveer 2025; Vitorino 2025 &
    Testing, ML-test, seguridad &
    Delimitación (fuera de foco) &
    Orthogonal al \AlignmentScore{} \\
    \textbf{E6} &
    Mahanra Rao 2025 &
    BERT + QoS en selección de servicios &
    S parcial &
    Recomendación $\neq$ alineación BIAN \\
    \bottomrule
  \end{tabular}
\end{table}


La Tabla~7 (Anexo~B) detalla objetivo, técnicas, métricas y limitaciones \textbf{paper por paper}.
